\section{Conclusiones}
\label{sec_6}
Se sintetizó un filtro pasabanda a partir de la plantilla de la Cátedra. La aproximación de Chebyshev de orden $2$, transformada a un pasabanda de orden $4$, coincide con la función de transferencia de referencia y ofrece $34.8 \, [dB]$ de atenuación en ambas bandas de rechazo. Se implementó con dos celdas Sallen-Key de componentes iguales ($k = 2.609$), un pasabajos en $1296 \, [Hz]$ y un pasaaltos en $771.5 \, [Hz]$, más un divisor con seguidor para ajustar la ganancia, utilizando tres de los cuatro amplificadores de un LM324. La simulación nominal reproduce la respuesta teórica: cumple el rechazo con más de $4.5 \, [dB]$ de margen, y el rizado resulta de $0.273 \, [dB]$ frente a los $0.25 \, [dB]$ especificados, diferencia que se debe a la cuantización de los valores normalizados.

El análisis de sensibilidad muestra la debilidad del diseño de componentes iguales. Aunque las sensibilidades de $\omega_p$ son de $-1/2$, las del ancho de banda llegan a $5.1$ respecto de los componentes y a $6.7$ respecto de $k$, porque el coeficiente de $s$ es proporcional a la diferencia $3 - k = 0.39$. Con componentes del $10 \, [\%]$ la frecuencia de los polos puede desplazarse entre $-17 \, [\%]$ y $+24 \, [\%]$, y el peor caso del ancho de banda es todavía más grave: en un sentido el $Q_p$ cae a $0.87$ y el rechazo no se cumple; en el otro, $k$ llega a $2.97$ y los polos pasan al semiplano derecho, de modo que el filtro oscila. La simulación de Montecarlo confirmó que esto no es un caso extremo improbable: el $8 \, [\%]$ de los circuitos con componentes del $10 \, [\%]$ resultó inestable.

El estudio de tolerancias mostró que el rechazo de $30 \, [dB]$ se cumple con holgura desde el $5 \, [\%]$, mientras que el rizado de $0.25 \, [dB]$ prácticamente no se alcanza ni con componentes del $1 \, [\%]$. Esta conclusión es propia de la aproximación de Chebyshev: al ser la solución óptima para el orden dado, su rizado nominal coincide con el límite y cualquier desvío lo excede. Para un armado real convendría diseñar con un rizado menor que el especificado, aprovechando el margen en las bandas de rechazo, o prever un ajuste de cada etapa.

La comparación con la topología de ganancia unitaria, con la misma función de transferencia, muestra que esta es claramente más robusta: la suma de las sensibilidades del ancho de banda baja de $13.3$ a $2$, no puede volverse inestable y, con componentes del $10 \, [\%]$, la mediana del rizado baja de $4.0$ a $0.9 \, [dB]$. Su costo es una mayor dispersión de valores ($C_1/C_2 = 26$). El amplificador real, por su parte, agrega unos $0.05 \, [dB]$ de rizado y limita la atenuación por encima de $40 \, [kHz]$, sin comprometer la especificación.

La automatización de las simulaciones con Python y \texttt{spicelib} fue la herramienta que hizo posibles estos análisis. A partir de los esquemáticos dibujados en LTspice, un único script ejecuta todas las simulaciones del trabajo (más de $3000$ corridas) en unos tres minutos, con resultados idénticos a los de las simulaciones manuales. Generar los valores desde Python permitió además calcular los polos de cada corrida y detectar la inestabilidad, que el análisis AC por sí solo no revela, y guardar los resultados permitió regenerar e inspeccionar las figuras sin volver a simular.

Por último, se diseñó un oscilador puente de Wien que, con $R = 8.25 \, [k\Omega]$ y $C = 1.5 \, [nF]$, oscila en $12042.6 \, [Hz]$, dentro de la tolerancia de $\pm 100 \, [Hz]$, con $2.85 \, [V]$ de amplitud y $3 \, [\%]$ de distorsión. El ajuste en simulación fue imprescindible, porque el retardo de fase del LM324 baja la frecuencia un $6 \, [\%]$ respecto de $1/(2\pi RC)$. También se verificó que la amplitud debe mantenerse por debajo del límite que impone el \textit{slew rate}, y que con componentes comerciales hace falta un preset para alcanzar la tolerancia pedida.
